\documentclass{article}
\usepackage{amsmath}
\usepackage{xcolor}
\usepackage{amssymb}
\usepackage{url}
\newcommand{\F}{\mathbb{F}}

\title{Writeup Regarding OPRF implementation}
\author{Diego F. Aranha \and Aron van Baarsen \and Adam Blatchley Hansen \and Kent Nielsen \and Peter Scholl}

\begin{document}
\maketitle

Alongside this write-up, the cleaned submission version of the implementations from the OPRF paper is attached.

\section{Structure of the implementation}
The implementation of the OPRF is an extension of the Swanky library \cite{swanky} written in Rust.
The implementation is part of a private GitHub Repository located at:
\begin{center}
    \url{https://github.com/dfaranha/Crypto-Dark-Matter-OPRF-minimal}
\end{center}
The implementation is available to reviewers in anonymized form, thus not publicly available yet. 
The choice of extending this publicly available library is due to several reasons, some of which are:
\begin{itemize}
    \item Number of building blocks needed already implemented (e.g, VOLE extension, OT extension, some field arithmetics).
    \item All implementations are generic, allowing for easy reuse of building blocks for different finite fields.
    \item The only library containing an implementation of the paper on high-degree ZK-proofs we utilize \cite{conversions}.
\end{itemize}
Hence, by extending this existing library, the implementation has been produced significantly faster than writing everything from scratch.

\section{Scope of our implementation}
The code contains the implementation of all protocols present in the paper, alongside additional elements needed to enable the protocol. 
The following is an exhaustive list of protocols introduced to the library:
\begin{itemize}
    \item Fully malicious OPRF.
    \item Malicious client OPRF.
    \item Semi-honest OPRF.
    \item One-sided authenticated Bits - Linear code approach.
    \item One-sided authenticated Bits - Cut \& Choose approach.
    \item Multiplication triples over $\F_3$.
    \item daBit protocol over $(\F_2, \F_3)$.
    \item Half authenticated Bits over $(\F_2, \F_3)$.
    \item Semi-honest modulus conversion bits.
    \item Input encoding and checks based on linear code.
\end{itemize}
Note that the Semi-honest variant is not included in the paper submission, and thus not proven secure. Usage of this variant requires caution.
The following is an exhaustive list of new parts introduced to the library to allow the above protocols to function.
\begin{itemize}
    \item Arithmetic in fields $(\F_3, \F_{3^{82}}, \F_{2^{256}})$.
    \item Endemic Oblivious Transfer based on Kyber.
    \item Kyber implementation.
    \item Optimizations to VOLE extension.
\end{itemize}
Note: Not all parts are novel from the paper/implementation, some are reused as described below.

\section{Reuse of Publicly Available Implementations \& Protocols}
The implementations of the following are either reused from existing code or protocols from publications.
\begin{itemize}
    \item Semi-honest modulus conversion bits: Implementation of \cite[Protocol 5.2]{C:DGHIKS21} by converting a OLE into an OT.
    \item Arithmetic on $(\F_3, \F_{3^{82}})$: The ternary field operations are implemented of algorithms from \cite{f3algo} and \cite{finitefieldmul}.
    \item Endemic Oblivious Transfer based on Kyber: The implementation is a Rust port of the C implementation from the library libOTe \cite{libOTe}.
    \item Kyber implementation: This is the implementation from \cite{Kyber}.
    \item Optimizations to VOLE extension: These optimizations include ideas from \cite{emp-zk} and new trivial changes to reduce the amount of work.
\end{itemize}

\section{Structure of code base}
The source code of the implementations can be found in the following locations.
\begin{itemize}
    \item OPRF Protocols - \texttt{oprf/src/*}
    \item Finite field arithmetic - \texttt{scuttlebutt/src/field/\{f3, f82t, f256b\}}
    \item Kyber implementation - \texttt{Kyber\_ot/*}
    \item Endemic OT - \texttt{ocelot/src/ot/endemic\_ot.rs}
    \item VOLE optimizations - \texttt{ocelot/src/svole/wykw/svole.rs}
\end{itemize}

\bibliographystyle{alpha}
\bibliography{code}
\end{document}
